\chapter{Requirement Analysis}
\label{ch:requirement-analysis}

The goal of this chapter is to produce a structured requirements specification for the application security data framework. The specification captures what the framework must do, at what quality level, and under what constraints, so that each requirement can later be traced to a concrete design decision and verified through automated tests.

%%%%%%%%%%%
%%% 2.1 %%%
%%%%%%%%%%%

\section{Requirements Methodology}
\label{sec:methodology}

\subsection{Requirements Sources}

Requirements were gathered from four complementary sources. First, a domain analysis of the application security landscape identified the data entities, tool categories, and integration patterns that any framework in this space must accommodate. Second, the author's professional experience designing and operating security data pipelines in large enterprises provided practical insight into real-world constraints and failure modes. Third, consultations with a professional network of security engineers, platform architects, and engineering managers validated the relevance and priority of candidate requirements. Fourth, the literature review conducted in \autoref{ch:literature-review} grounded the requirements in established data engineering and data governance principles.

\subsection{Data Contracts as Specification Format}
\label{sec:data-contracts}

This thesis specifies requirements using \emph{data contracts}, a concept drawn from modern data engineering practice. A data contract is a formal, structured agreement that defines the schema, semantics, quality rules, and service level expectations for a data asset \cite{Jones2023}. The DAMA Data Management Body of Knowledge similarly advocates for explicit specification of data quality dimensions and governance metadata as part of any data management initiative \cite{DAMA2017}.

Data contracts are well suited for specifying data pipeline requirements because they capture precisely the information that must be agreed upon between producers and consumers of data: what inputs are expected, what outputs are produced, which quality rules must be in place, and what service level agreements apply. Unlike free-form requirement documents, data contracts are structured and machine-readable, which enables automated validation and traceability.

Each requirement in this thesis is expressed as a data contract containing the fields described in \autoref{tab:contract-fields}. \autoref{code:contract-example} shows an example contract for a representative requirement.

\begin{table}[ht]
\centering
\caption{Data contract template fields.}
\label{tab:contract-fields}
\begin{tabular}{l p{9cm}}
\hline
\textbf{Field} & \textbf{Description} \\
\hline
\texttt{id} & Unique identifier following the \texttt{REQ-AREA-NNN} pattern. \\
\texttt{name} & Short, descriptive name of the requirement. \\
\texttt{domain} & Functional domain the requirement belongs to (see \autoref{tab:req-domains}). \\
\texttt{priority} & MoSCoW priority level (see \autoref{sec:moscow}). \\
\texttt{persona} & Primary stakeholder persona the requirement serves. \\
\texttt{description} & Prose statement of what the requirement demands. \\
\texttt{inputs} & Data sources or artifacts consumed. \\
\texttt{outputs} & Data assets or artifacts produced. \\
\texttt{quality\_rules} & Verifiable conditions that must hold for the requirement to be satisfied. \\
\texttt{sla} & Service-level targets such as latency, freshness, or availability. \\
\hline
\end{tabular}
\end{table}

\begin{code}{}{Example data contract for a security finding ingestion requirement.}{code:contract-example}
id: REQ-ING-001
name: Ingest Security Findings
domain: ING
priority: Must Have
persona: Security Expert
description: >
  The framework shall ingest security findings from
  source tool APIs and persist them for downstream
  processing.
inputs:
  - source_tool_api_responses
outputs:
  - raw_security_findings
quality_rules:
  - all records must contain a non-null finding identifier
  - severity must be one of critical, high, medium, low
sla:
  freshness: 1 hour
  availability: 99.5%
\end{code}

\subsection{Test-to-Requirement Linking}

The framework employs two complementary verification mechanisms aligned with Databricks development patterns.

\subsubsection{Unit and Integration Tests}

Python modules containing connector logic and transformation functions are tested using pytest with custom requirement markers. \autoref{code:req-marker} demonstrates how tests reference requirement IDs.

\begin{code}{Python}{Requirement Marker Example}{code:req-marker}
@pytest.mark.requirement("REQ-ING-002")
def test_sonarqube_pagination():
    """Validate pagination handling for large result sets."""
    ...
\end{code}

\subsubsection{Data Quality Expectations}

Data quality requirements are verified using Delta Live Tables expectations, which execute during pipeline runs. \autoref{code:dlt-expectation} shows how expectations map to requirements via code comments.

\begin{code}{Python}{DLT Expectation Example}{code:dlt-expectation}
@dlt.table
@dlt.expect_or_fail("valid_severity", "severity IN ('critical', 'high', 'medium', 'low')")
@dlt.expect("has_finding_id", "finding_id IS NOT NULL")
def silver_findings():
    # REQ-NFR-001: Data Quality Validation
    return transform_findings(dlt.read("bronze_raw"))
\end{code}

Expectation results are captured in the pipeline event log and reflected in requirement status.

\subsection{Traceability Generation}

A generator script produces the traceability matrix by:
\begin{enumerate}
    \item Scanning test files for \texttt{@pytest.mark.requirement} markers
    \item Parsing DLT pipeline code for expectation comments referencing requirement IDs
    \item Collecting pytest results and DLT pipeline event logs
    \item Merging with requirement descriptions from specification
\end{enumerate}



Appsec vendors, products and their capabilities, APIs and their maturity/completeness

Capture data formats and integration patterns of the tools


%%%%%%%%%%%
%%% 2.2 %%%
%%%%%%%%%%%

\section{Personas}
\label{sec:personas}

The framework serves several distinct stakeholder groups, each with different data needs and expectations. Understanding these personas is essential for deriving the requirements that the framework must satisfy. The following subsections characterize the four primary persona categories whose needs shape the framework.

\subsection{Application Owners}

Application owners are business stakeholders responsible for the overall risk posture of one or more applications within the enterprise portfolio. They need aggregated security posture visibility, including trend data over time and composite risk metrics that communicate the security state of their applications at a glance. Rather than examining individual findings, application owners interact with dashboards that present overall health indicators, remediation progress, and compliance status. Their need for technical detail is limited; they require information in the form of actionable business terms. These needs directly drive the framework's aggregation and serving layer requirements, particularly the computation of application-level risk scores and trend visualizations.

\subsection{Application Developers}

Application developers are the engineers who, in the context of our framework, receive security findings and are responsible for their remediation. They need clear, actionable vulnerability information enriched with code-level context such as the affected file, line number, and remediation guidance. Developers have a low tolerance for false positives and duplicated noise, as excessive or inaccurate findings erode trust in the security tooling and reduce remediation flow. Findings must be presented in a format that integrates naturally into their existing workflows, whether through issue trackers or \gls{cicd} pipeline feedback. These needs drive the framework's normalization, de-duplication, and data quality requirements, ensuring that only validated, contextually enriched findings reach developers.

\subsection{Security Experts}

Security experts are the security architects and security engineers who support application teams, auditors responsible for assessing security compliance on application level, or security analysts seeking context for a security event or incident. They require access to detailed findings, complete audit trails, and the ability to drill down into raw data when performing triage or validating tool output. Security engineers define and maintain security policies, evaluate tool effectiveness, and coordinate remediation priorities across teams. Compliance auditors, a closely related sub-group, need evidence of security requirement coverage, remediation timelines, and historical records that demonstrate adherence to organizational and regulatory standards. These needs drive the framework's data quality, governance, and granular access control requirements, as well as the preservation of full data lineage from source to consumption.

\subsection{Development and Security Leadership}

Development and security leadership, including \glspl{ciso}, engineering managers, and vice presidents, require quantified risk metrics to inform strategic decisions and resource allocation. Key indicators include \gls{mttr}, \gls{sla} compliance rates, vulnerability density trends, and team-level or portfolio-level comparisons that reveal where remediation investment is most needed. These stakeholders consume information through executive dashboards designed for rapid interpretation rather than detailed investigation. Their decisions directly affect staffing, tooling budgets, and organizational security strategy. These needs drive the framework's gold-layer aggregation and metric computation requirements, demanding reliable, well-governed data pipelines that produce trustworthy summary metrics.

%%%%%%%%%%%
%%% 2.3 %%%
%%%%%%%%%%%

\section{Data Sources}

The framework must support ingestion of data from a diverse set of source systems to serve all personas identified in \autoref{sec:personas}. As discussed in the introduction, security data in a typical enterprise is fragmented across dozens of tools, each with its own \gls{api}, data format, and update frequency. This fragmentation prevents stakeholders from obtaining a unified view of application security posture.

This section analyzes each source category in four perspectives: what data the source provides and why it matters, how that data can be accessed programmatically (\gls{api} analysis), what formats and protocols are involved, and what integration challenges the framework must address. The analysis proceeds from business context (application inventory), through development infrastructure (\gls{scm}, issue trackers, \gls{cicd}), to runtime environments and, in later subsections, to the security testing tools themselves.

\subsection{Application Inventory}

The application inventory is the foundational data source that establishes the business context for all security findings. Without a reliable catalog of business applications and their mapping to technical assets, a vulnerability discovered in a code repository or a running service has no business meaning: it cannot be attributed to an application owner, assessed for business impact, or prioritized against organizational risk criteria.

Enterprise organizations typically maintain their application inventories in configuration management databases (\glspl{cmdb}) such as ServiceNow, in custom-built internal catalogs, or in a combination of both. These registries capture metadata about each business application, including its name, owning team, lifecycle status, and business unit. Critically, they also hold business impact assessment data: criticality ratings (commonly structured as tiers, such as Tier~1 for mission-critical applications, Tier~2 for important supporting systems, and Tier~3 for non-critical internal tools), data classification levels indicating the sensitivity of information processed, and regulatory scope flags that identify which compliance frameworks apply to a given application.

TODO API vs ServiceNow application

\subsection{Tools for Software Development and Delivery}

important data sources for inventory

start with SCM and issue tracking tools

Github, Gitlab, Bitbucket, Azure DevOps

Jira, Azure DevOps, Github and Gitlab issues

CI/CD pipelines - Jenkins, Github Actions, Azure DevOps

\subsection{Runtime Infrastructure Platforms and Security Tools}

We need infra portrayed from app angle - what does the app run on, how is it secure, whats the security context

AWS - EC2, serverless
Azure
GCP

Also how secure is it (infra VMDR like Qualys, Wiz...) can impact risk of app vulns if they run on vulnerable infra too

How to get data from Kubernetes, Docker even if on prem

\subsection{Security Testing Tools}

Primary vulnerability discovery tools (secret scanning, , , ...)

\subsubsection{Secret Scanning Tools}

\subsubsection{Static Application Security Testing (SAST) Tools}

\gls{sast}

\subsubsection{Software Composition Analysis (SCA) Tools}

\gls{sca}

\subsubsection{Dynamic Application Security Testing (DAST) Tools}

\gls{dast}

\subsubsection{Penetration Testing}

Need a way to upload manual findings, e.g. PDF report from external pentesters


%%%%%%%%%%%
%%% 2.4 %%%
%%%%%%%%%%%

\section{Data Transformations}

\subsection{Data Ingestion and Validation}

Filtering?

\subsection{Data Normalization}

Standard format for each object and finding type

\subsection{Finding Deduplication}

When 2 tools find the same thing, link it, don't create 2 issues

\subsection{Business Application Mapping}

\subsection{Triage}


%%%%%%%%%%%
%%% 2.5 %%%
%%%%%%%%%%%

\section{Data Consumers}

%%% 2.5.1 %%%
\subsection{Remediation Tracking Systems}

Jira, ServiceNow

%%% 2.5.2 %%%
\subsection{Application Security Posture Management (ASPM)}

Can pull from data platform already cleaned data instead of from each source separately

%%% 2.5.3 %%%
\subsection{SOAR and SIEM}


%%%%%%%%%%%
%%% 2.6 %%%
%%%%%%%%%%%

\section{Non-Functional Requirements}

\subsection{Data Quality}

Data validations and transformations on input are critical

\subsection{Extensibility}

Important feature of the whole framework, especially on data source integrations and ML models

\subsection{Scalability}

it's intended for large enterprises, 10k's of repos +

\subsection{Performance}

Some use cases require low latency

\subsection{Security}

Obviously, with this use case

\subsection{Maintainability}

SaaS preferred over manual maintenance - Databricks is ideal

All source code versioned in git and following best practices



%%%%%%%%%%%
%%% 2.8 %%%
%%%%%%%%%%%

\section{Requirement Specification}

\subsection{Data Ingestion}

\subsection{Data Processing}

\subsection{Data Serving}

\subsection{Non-Functional Requirements}
